\chapter{Related Work}
\label{sec:lavori_correlati}

Negli ultimi anni la sicurezza dei sistemi Industrial Control Systems (ICS) e Supervisory Control and Data Acquisition (SCADA) è diventata una priorità cruciale, spingendo la ricerca verso lo sviluppo di strumenti e metodologie per la rilevazione e la mitigazione delle minacce informatiche. Diversi lavori hanno contribuito alla definizione di ambienti di simulazione e modelli di difesa integrati, con l’obiettivo di testare e migliorare la resilienza dei sistemi critici. Tra questi, alcuni studi iniziali hanno proposto l’implementazione e la simulazione di servizi di sicurezza in infrastrutture SCADA reali e virtualizzate, evidenziando l’importanza di valutare il comportamento dei sistemi sotto attacco e la risposta dei meccanismi di difesa in scenari realistici \cite{simulationImplementationSecurityServices}.

Ulteriori ricerche hanno introdotto modelli multilivello di protezione, come il sistema a triplo livello per la sicurezza SCADA in reti elettriche, basato sull’interazione coordinata tra livelli fisici, logici e applicativi \cite{tripleLayerScadaSecurityOfElectricUtility}. Questo approccio, pensato per le infrastrutture critiche, ha aperto la strada a nuove architetture difensive che combinano monitoraggio continuo, gestione delle anomalie e risposta automatica agli incidenti. Parallelamente, altri studi si sono concentrati sull’architettura di sicurezza informatica per la smart grid, sottolineando la necessità di una progettazione modulare che consenta la protezione di ogni componente del sistema \cite{architectingCyberSecSmartGrid}.

Un contributo significativo è rappresentato dal progetto PRECINCT, che ha proposto un ecosistema di cooperazione pubblico-privato per migliorare la resilienza delle infrastrutture critiche urbane attraverso la condivisione dei dati e la correlazione degli eventi di sicurezza \cite{precinct}. Tale approccio è stato affiancato da sistemi specifici per la rilevazione delle intrusioni, come DIDEROT IDS, un framework per la difesa in ambienti industriali basato su tecniche di Machine Learning e su moduli di analisi del traffico \cite{diderotIds}.

L’analisi dei pacchetti e l’identificazione dei dispositivi industriali tramite il payload dei messaggi rappresentano un’altra direzione di ricerca rilevante, volta a migliorare la consapevolezza del sistema e la profilazione dei dispositivi connessi \cite{identificationPayloadIndustrialDevices}. Parallelamente, l’esplorazione della diversità nelle architetture SCADA ha permesso di comprendere meglio le vulnerabilità e le possibili superfici d’attacco \cite{scadaDiversityExploration}.

Sono stati inoltre sviluppati approcci basati su regole per il monitoraggio delle reti SCADA \cite{ruleBasedScadaNetwork} e strumenti come Suricata, che, attraverso l’estensione delle proprie funzionalità, si sono dimostrati efficaci nel rafforzare la sicurezza informatica in contesti ICS \cite{suricataEnhancingIdsCyberSecurity}. Altri lavori hanno applicato metodi di Machine Learning supervisionato per la rilevazione delle intrusioni, mostrando come la classificazione automatica degli eventi possa incrementare la capacità di identificare comportamenti anomali \cite{supervisedLearningIdsScada}.

A questi studi si aggiungono ricerche che propongono un approccio fisico-basato alla rilevazione delle intrusioni, combinando metriche di rete e misure fisiche dei processi per migliorare l’accuratezza degli IDS e ridurre i falsi positivi nei sistemi SCADA \cite{alasiri202034}.

In aggiunta, framework specifici hanno integrato la semantica dei processi industriali per migliorare la precisione nella rilevazione delle minacce e ridurre i falsi positivi \cite{idsScadaFrameworkIncorporateProcessSemantics, semanticSecurityAnalysisScadaNetwork}. Tali approcci semantici hanno consentito di analizzare le relazioni tra dati di processo e attività di rete, fornendo una visione più completa della sicurezza operativa.

Recentemente, la ricerca si è spostata verso l’uso di tecniche di difesa attiva e modelli basati su grafi, capaci di rappresentare le relazioni tra componenti di rete e potenziali vettori d’attacco \cite{graphBasedActiveCyberDefense}. In particolare, l’applicazione di Graph Neural Network (GNN) ha mostrato risultati promettenti nella classificazione e nella previsione di eventi di sicurezza nei sistemi elettrici industriali \cite{graphNeuralNetworkIdsElectrical}.

Anche il tema della valutazione del rischio nei PLC (Programmable Logic Controller) è stato affrontato con approcci quantitativi e metodologie per l’identificazione delle vulnerabilità \cite{plcCybersecuirtyRisk}. Infine, sono emersi filoni di ricerca che combinano l’intelligenza artificiale con la sicurezza informatica in domini specifici, come l’agricoltura intelligente \cite{cyberSecurityAgricoltureMachineLearning} e la protezione dei sistemi SCADA mediante tecniche AI avanzate \cite{aiScada}, confermando la tendenza verso l’adozione di strumenti intelligenti per la difesa delle infrastrutture critiche.